\documentclass[10pt,a4paper]{amsart}
\usepackage[margin=19mm]{geometry}
\usepackage{amsmath,amssymb,mathtools}
\usepackage[scr=rsfs,cal=euler]{mathalfa}
\usepackage{xfrac}
\usepackage[unicode,hidelinks,bookmarks=false]{hyperref}
\newcommand{\ie}{i.e.\ }
\newcommand{\cf}{cf.\ }
\newcommand{\resp}{respectively\ }
\newcommand{\Subsection}{\subsection}
\providecommand{\nobreakdash}{\nobreak\hbox{-}}
\setlength{\emergencystretch}{2em}
\multlinegap0pt
\newcommand{\fl}[1]{\left\lfloor #1\right\rfloor}
\newcommand{\md}{\mathbin{\mathrm{mod}}}
\theoremstyle{plain}
\newtheorem{theorem}{Theorem}
\newtheorem{lemma}[theorem]{Lemma}
\newtheorem{proposition}[theorem]{Proposition}
\newtheorem{corollary}[theorem]{Corollary}
\theoremstyle{remark}
\newtheorem{remark}[theorem]{Remark}
\title[Straight-line programs for the solutions of Pell's equation]{Straight-line programs for the solutions of Pell's equation}
\author{Bogdan Dumitru, Mihai Prunescu}
\date{Source: arXiv:2609.28649v1 (CC BY 4.0)}
\begin{document}
\maketitle

\noindent\emph{Source and licence.} The delivered work is the article shown in the title above, version arXiv:2609.28649v1 (CC BY 4.0), distributed under the Creative Commons Attribution 4.0 International licence, as recorded in the arXiv licence metadata for this exact version and shown on the abstract page saved with this item. The full licence notice, which carries the licence URL and the address of that abstract page, is the bundled file \texttt{licenses/arxiv-2609.28649-CC-BY-4.0.txt}.
\par\smallskip

\noindent\textit{Editorial scope.} Counted unit: the article's Theorem (source label thm:general-base) (source0.tex lines 1665--1673, source label \texttt{thm:general-base}): ``Least admissible common base The least admissible common base for and is b 0 2X 1(X 1 1)-1. Both formulas hold for all and all integers . For every integer , the pair fails at .''. It is delivered together with the article's complete original proof (source0.tex lines 1675--1757, one proof environment adjacent to the statement, 83 lines), copied line for line from the article's own text. Required context retained uncounted: eq:general-x (lines 1617--1622); lem:tail (lines 1639--1650). Editorial formatting only: an AMS \texttt{amsart} preamble that restates the article's own macro definitions and its own theorem declarations, and the theorem environments the retained lines use; the article's own class file is neither shipped nor input; the retained environments are renumbered sequentially in order of appearance, whereas the article prints them with its own numbering; the article's equation numbering is not reproduced and every cross-reference inside the retained text resolves to the wrapper's numbering. No citation occurs inside any retained line, so the wrapper carries no bibliography; All retained bodies are exact copies of the bundled \texttt{sources/211606/source0.tex}, so no omission receipt applies. No asset-inclusion command occurs inside any retained span and the package ships no image asset.


\section*{Retained context: eq:general-x (source0.tex lines 1617--1622)}

\begin{align}
  X_n&=\fl{\frac{b^{n^2+2n}-X_1b^{n^2+n}}
                  {b^{2n}-2X_1b^n+1}}\md b^n,\label{eq:general-x}\\
  Y_n&=\fl{\frac{Y_1b^{n^2+n}}
                  {b^{2n}-2X_1b^n+1}}\md b^n.\label{eq:general-y}
\end{align}


\section*{Retained context: lem:tail (source0.tex lines 1639--1650)}

\begin{lemma}\label{lem:tail}
Let $\Lambda=b^n>2X_1$, and define
\[
  T_n=\sum_{i\ge1}X_{n+i}\Lambda^{-i},\qquad
  U_n=\sum_{i\ge1}Y_{n+i}\Lambda^{-i}.
\]
The floor in \eqref{eq:general-x} equals
$\sum_{j=0}^n X_j\Lambda^{n-j}+\fl{T_n}$.
If $X_n<\Lambda$ and $T_n<1$, formula \eqref{eq:general-x}
returns $X_n$.
The corresponding statements hold for $Y_n$ and $U_n$.
\end{lemma}


\section{The counted unit: Theorem (source label thm:general-base) and its complete original proof}

\begin{theorem}[Least admissible common base]\label{thm:general-base}
The least admissible common base for \eqref{eq:general-x} and
\eqref{eq:general-y} is
\begin{equation}\label{eq:base}
  b_0=2X_1(X_1+1)-1.
\end{equation}
Both formulas hold for all $n\ge1$ and all integers $b\ge b_0$.
For every integer $2X_1\le b<b_0$, the pair fails at $n=1$.
\end{theorem}

\begin{proof}
Write $a=X_1$ and $y=Y_1$ in this proof. Then $a\ge2$,
$1\le y<a$, and $\gcd(a,y)=1$.

First consider $b=b_0$ and $n=1$. Summing the generating-function
tail gives
\[
  T_1=\frac{b(2a^2-1)-a}{b^2-2ab+1}.
\]
For $b>2a$, the inequality $T_1<1$ is equivalent to
\begin{equation}\label{eq:base-quadratic}
  f(b)=b^2-(2a^2+2a-1)b+(a+1)>0.
\end{equation}
At $b_0$, one has $f(b_0)=a+1>0$.
Since $a<b_0$ and $U_1<T_1$, Lemma~\ref{lem:tail} proves both
formulas at $n=1$.

For $n\ge2$, the recurrence implies $X_m<(2a)^m$ for $m\ge1$.
The inequality
\[
  b_0^2-(2a)^3-2a=4a^4-6a+1>0
\]
starts an induction giving
\[
  b_0^n>(2a)^{n+1}+2a\qquad(n\ge2).
\]
To pass to the next value of $n$, multiply by $b_0>2a$.
Consequently
\[
  T_n<\sum_{i\ge1}(2a)^{n+i}b_0^{-ni}
      =\frac{(2a)^{n+1}}{b_0^n-2a}<1.
\]
Also $X_n<(2a)^n<b_0^n$ and $U_n<T_n$.
Lemma~\ref{lem:tail} proves both extraction formulas.
Increasing $b$ decreases the tails and increases the modulus, proving
sufficiency for $b\ge b_0$.

It remains to exclude smaller bases. At $b=2a$, the denominator for
$n=1$ is one, and the $X$-formula returns zero modulo $b$.
Suppose now $b\ge2a+1$, and set $s_b=b-2a$ and $D_b=bs_b+1$.
The two tails at $n=1$ are
\[
  T=\frac{b(2a^2-1)-a}{D_b},\qquad
  U=\frac{y(2ab-1)}{D_b}.
\]
The floors before reduction modulo $b$ are
$b+a+\fl T$ and $y+\fl U$.
If both residues are correct, then
$\fl T=kb$ and $\fl U=lb$ for nonnegative integers $k,l$.
Put $H_b=2a^2-1-kD_b$ and $J_b=2ay-lD_b$.
The floor inequalities give
\[
  \frac ab\le H_b<s_b+\frac{a+1}{b},\qquad
  \frac yb\le J_b<s_b+\frac{y+1}{b}.
\]
Thus $1\le H_b,J_b\le s_b$.
On the other hand,
\[
  yH_b-aJ_b+y=(al-yk)D_b,
\]
and the left side lies strictly between $-D_b$ and $D_b$, because
\[
  -D_b<2y-as_b\le yH_b-aJ_b+y
       \le y(s_b+1)-a<as_b<D_b.
\]
Hence $al=yk$. Coprimality gives $a\mid k$, whereas
\[
  0\le k\le T/b<2a^2/D_b<a
\]
because $D_b>2a$. Therefore $k=l=0$, and simultaneous correctness
requires $T<1$.

The quadratic in \eqref{eq:base-quadratic} has one root in $(0,1)$
and the other in $(b_0-1,b_0)$: indeed
\[
  f(0)=a+1>0,\qquad
  f(1)=f(b_0-1)=3-a-2a^2<0,\qquad f(b_0)=a+1>0.
\]
For integers $b>2a$, the condition $T<1$ is therefore equivalent
to $b\ge b_0$.
For the remaining bases $2\le b<2a$, the denominator at $n=1$ is
$b(b-2a)+1<0$, so these bases are outside the stated domain.
\end{proof}

\end{document}
